\documentclass{article}
\usepackage[T1]{fontenc}
\usepackage{lmodern}
\usepackage{graphicx}
\usepackage{amsmath,amssymb}
\usepackage[a4paper,margin=2cm]{geometry}
\usepackage[absolute,overlay]{textpos}
\setlength{\TPHorizModule}{1mm}
\setlength{\TPVertModule}{1mm}
\usepackage[indonesian]{babel}
\usepackage{hyperref}
\setlength{\emergencystretch}{2em}
% unit: Halaman 4

\begin{document}

% === Halaman 4 (python/native) ===

4

Properti Algoritma RSA

1.   p dan q bilangan prima 


(rahasia)

2.   n = p  q 




(tidak rahasia)

3.   (n) = (p – 1)(q – 1) 


(rahasia)

4.   e     (kunci enkripsi) 


(tidak rahasia, disebut kunci publik)

     Syarat: PBB(e, (n)) = 1   , PBB = pembagi bersama terbessar = gcd

5.  d     (kunci dekripsi) 


(rahasia, disebut kunci privat)

     d    dihitung dari d  e-1 mod ((n) )

6.  m    (plainteks) 



(rahasia)

7.  c    (cipherteks) 



(tidak rahasia)

\end{document}
